\documentclass[11pt,a4paper]{article}
\usepackage[margin=1in]{geometry}
\usepackage{amsmath,amssymb,amsthm}
\usepackage{algorithm}
\usepackage{algpseudocode}
\usepackage{booktabs}
\usepackage{hyperref}

\theoremstyle{definition}
\newtheorem{definition}{Definition}[section]
\newtheorem{theorem}{Theorem}[section]
\newtheorem{lemma}[theorem]{Lemma}
\newtheorem{remark}{Remark}[section]

\title{\textbf{Rigorous Mathematical Specification of Retrieval-Enhanced Transformer Chunks (RETRO)}}
\author{Team Scaibu \\ \small Email: \texttt{jaydeep.9664920749@gmail.com} \\ \small Architecture and Core Engine Team}
\date{October 2026}

\begin{document}
\maketitle

\begin{abstract}
Retrieval-Enhanced Transformers (RETRO) augment dense language modeling with an external retrieval database. Sequences are partitioned into chunks; each chunk retrieves $k$ nearest neighbors whose representations cross-attend with input tokens, decoupling memorization from parameter scale.
\end{abstract}

\begin{thebibliography}{9}
\bibitem{borgeaud2022} S.~Borgeaud et al., ``Improving language models by retrieving from trillions of tokens,'' in \emph{ICML}, 2022.
\end{thebibliography}

\end{document}
